\begin{proofbox}
\textbf{\textcolor{CProof}{证明目标与条件}}

设 $a>b>0$，$P(x_0,y_0)$ 在
$E:\dfrac{x^2}{a^2}+\dfrac{y^2}{b^2}=1$ 上，即
$\dfrac{x_0^2}{a^2}+\dfrac{y_0^2}{b^2}=1$。
证明 $P$ 处切线为
\[
\boxed{\frac{x_0x}{a^2}+\frac{y_0y}{b^2}=1.}
\]

\medskip
\textbf{\textcolor{CProof}{方法一｜最快代数证明：柯西不等式（推荐速证）}}

\textbf{第一步：} 对椭圆上任意点 $Q(x,y)$，应用柯西不等式：
\[
\begin{aligned}
\frac{x_0x}{a^2}+\frac{y_0y}{b^2}
&=\frac{x_0}{a}\frac{x}{a}+\frac{y_0}{b}\frac{y}{b}\\
&\le\sqrt{\frac{x_0^2}{a^2}+\frac{y_0^2}{b^2}}
\sqrt{\frac{x^2}{a^2}+\frac{y^2}{b^2}}=1.
\end{aligned}
\]
\textbf{第二步：} 等号成立当且仅当
$\left(\dfrac{x}{a},\dfrac{y}{b}\right)
=\left(\dfrac{x_0}{a},\dfrac{y_0}{b}\right)$，即 $Q=P$。
故直线 $\dfrac{x_0x}{a^2}+\dfrac{y_0y}{b^2}=1$
仅在 $P$ 处与椭圆相遇，且椭圆完全位于其一侧，必为切线。

\textbf{优势：} 两步完成；不分类讨论 $y_0=0$，适合需要迅速核验证明的题目。

\medskip
\textbf{\textcolor{CProof}{方法二｜几何证明：椭圆伸缩为单位圆（最直观）}}

\textbf{第一步：} 作坐标伸缩
\[
 u=x/a,\quad v=y/b,\quad u_0=x_0/a,\quad v_0=y_0/b.
\]
于是椭圆变成单位圆 $u^2+v^2=1$，$P$ 变为 $P'(u_0,v_0)$。

\textbf{第二步：} 圆在 $P'$ 处的切线垂直于半径 $OP'$，
故切线上点 $(u,v)$ 满足
\[
 u_0(u-u_0)+v_0(v-v_0)=0.
\]
展开并用 $u_0^2+v_0^2=1$ 得
\[
 u_0u+v_0v=1.
\]

\textbf{第三步：} 代回原坐标得到
\[
 \frac{x_0x}{a^2}+\frac{y_0y}{b^2}=1.
\]
坐标伸缩为可逆线性变换，直线仍映为直线；圆切线与圆的唯一公共点对应椭圆切线与椭圆的唯一公共点，因而相切关系保留。

\medskip
\textbf{\textcolor{CProof}{方法三｜法向量：已知垂直方向就能写直线}}

\textbf{第一步：} 把椭圆写作
\[
F(x,y)=\frac{x^2}{a^2}+\frac{y^2}{b^2}-1=0.
\]
等值曲线的法向量为 $\nabla F$，在 $P$ 处
\[
 \nabla F(P)=\left(\frac{2x_0}{a^2},\frac{2y_0}{b^2}\right).
\]
这里的法向量也可由方法二的圆切线经过坐标伸缩直接得到。

\textbf{第二步：} 直线过 $P$，方向向量与该法向量垂直：
\[
 \frac{x_0}{a^2}(x-x_0)+\frac{y_0}{b^2}(y-y_0)=0.
\]
\textbf{第三步：} 移项并利用 $P$ 在椭圆上的条件，
\[
 \frac{x_0x}{a^2}+\frac{y_0y}{b^2}
 =\frac{x_0^2}{a^2}+\frac{y_0^2}{b^2}=1.
\]
\textbf{关键认识：} 法向量不是一般意义下的半径向量 $(x_0,y_0)$。

\medskip
\textbf{\textcolor{CProof}{方法四｜隐函数求导：熟悉的切线斜率法}}

\textbf{第一步：} 对 $\dfrac{x^2}{a^2}+\dfrac{y^2}{b^2}=1$ 两边关于 $x$ 求导：
\[
 \frac{2x}{a^2}+\frac{2y}{b^2}y'=0.
\]
\textbf{第二步：} 当 $y_0\ne0$ 时，切线斜率
\[
 k=y'(x_0)=-\frac{b^2x_0}{a^2y_0}.
\]
故
\[
\begin{aligned}
 y-y_0&=-\frac{b^2x_0}{a^2y_0}(x-x_0),\\
 b^2x_0x+a^2y_0y&=b^2x_0^2+a^2y_0^2=a^2b^2,
\end{aligned}
\]
两边除以 $a^2b^2$ 得所证。

\textbf{特殊情况：} 若 $y_0=0$，因 $P$ 在椭圆上，$x_0=\pm a$；
由椭圆在左右端点处的竖直切线知 $x=\pm a$，也正符合统一切线方程。

\medskip
\textbf{\textcolor{CProof}{方法五｜保留原思路：联立点斜式，利用重根}}

本方法先取 $y_0\ne0$，在此情况下切线不是竖直直线。

\textbf{第一步：} 设过 $P$ 的直线为 $y=y_0+k(x-x_0)$，代入椭圆：
\[
 \frac{x^2}{a^2}+\frac{[y_0+k(x-x_0)]^2}{b^2}=1.
\]
\textbf{第二步：} 用 $\dfrac{x_0^2}{a^2}+\dfrac{y_0^2}{b^2}=1$ 消去常数，整理：
\[
 \frac{x^2-x_0^2}{a^2}
 +\frac{2ky_0(x-x_0)}{b^2}
 +\frac{k^2(x-x_0)^2}{b^2}=0.
\]
提取 $x-x_0$：
\[
 (x-x_0)\left[
 \frac{x+x_0}{a^2}+\frac{2ky_0}{b^2}
 +\frac{k^2(x-x_0)}{b^2}\right]=0.
\]
\textbf{第三步：} 相切时 $x_0$ 是二重根，方括号在 $x=x_0$ 时必须为 $0$：
\[
 \frac{2x_0}{a^2}+\frac{2ky_0}{b^2}=0
 \quad\Longrightarrow\quad k=-\frac{b^2x_0}{a^2y_0}.
\]
\textbf{第四步：} 把 $k$ 代回点斜式：
\[
\begin{aligned}
 y-y_0&=-\frac{b^2x_0}{a^2y_0}(x-x_0),\\
 b^2x_0x+a^2y_0y
 &=b^2x_0^2+a^2y_0^2=a^2b^2,
\end{aligned}
\]
从而 $x_0x/a^2+y_0y/b^2=1$。

\textbf{补全端点：} $y_0=0$ 时单独取 $x=\pm a$。
\textbf{注：} 上述因式分解是对原证明的完整保留与压缩；“重根”发生在有限斜率直线对应的 $x$ 方程中，不能据此漏掉竖直切线。

\medskip
\textbf{\textcolor{CProof}{方法六｜先判定直线相切，再反求切点（拓展）}}

\textbf{第一步：} 对非竖直直线 $y=kx+m$，代入椭圆并乘以 $a^2b^2$：
\[
 (b^2+a^2k^2)x^2+2a^2kmx+a^2(m^2-b^2)=0.
\]
\textbf{第二步：} 判别式为
\[
\begin{aligned}
 \Delta
 &=4a^4k^2m^2-4a^2(b^2+a^2k^2)(m^2-b^2)\\
 &=4a^2b^2(a^2k^2+b^2-m^2).
\end{aligned}
\]
相切 $\iff\Delta=0\iff\boxed{m^2=a^2k^2+b^2}$，特别地 $m\ne0$。

\textbf{第三步：} 重根即切点横坐标：
\[
 x_0=-\frac{a^2km}{b^2+a^2k^2}
 =-\frac{a^2k}{m}.
\]
将其代回直线得
\[
 y_0=kx_0+m=\frac{m^2-a^2k^2}{m}=\frac{b^2}{m}.
\]
\textbf{第四步：} 由切点坐标反过来验证切点公式：
\[
 \frac{x_0x}{a^2}+\frac{y_0y}{b^2}
 =-\frac{kx}{m}+\frac{y}{m}=1
 \iff y=kx+m.
\]
\textbf{价值：} 把“已知切点求切线”和“已知斜率求切线”两类题连起来。
竖直切线 $x=\pm a$ 另行处理。

\medskip
\textbf{\textcolor{CProof}{补充校验｜平方和恒等式：确认真的只有一个公共点}}

若点 $Q(x,y)$ 同时满足椭圆和上述候选切线，则
\[
\begin{aligned}
&\frac{(x-x_0)^2}{a^2}+\frac{(y-y_0)^2}{b^2}\\
&=\left(\frac{x^2}{a^2}+\frac{y^2}{b^2}\right)
 -2\left(\frac{x_0x}{a^2}+\frac{y_0y}{b^2}\right)
 +\left(\frac{x_0^2}{a^2}+\frac{y_0^2}{b^2}\right)\\
&=1-2+1=0.
\end{aligned}
\]
两个平方均非负，只能同时为零，所以 $Q=P$。这也直接验证了端点情况。

\medskip
\textbf{\textcolor{CProof}{结论回扣}}

\textbf{背诵优先级：} 已知切点先记“平方变乘积”；讲清几何本质选“椭圆压成圆”；
需要最快严格验证用“柯西不等式”；求含 $k,m$ 的切线选“判别式 $\Delta=0$”。
\end{proofbox}
